\chapter[Symbolic representation analysis and barycenter estimation]{Symbolic Representation Analysis and Barycenter Estimation for the Assessment of dysexecutive syndromes}
\label{sec:chapter_6}

\begin{chapterabstract}

    % Alphabet c'est plus approprié parce que ça represented des objets plus petis (de semantiques plus equivoques) que le terme de vocabulaire
    % This chapter builds on the symbolic representation of participants’ egocentric vision introduced in the previous chapter to develop novel mathematical analyses of task execution patterns. 

    \textbf{Key points:}
    \begin{enumerate}
        \item This chapter investigates the use of symbolic representation of egocentric vision to study inter-individual variability and differences in task execution between control participants and patients with dysexecutive syndromes (TBI and RIL).

        \item We exploit the discriminative structure of symbols’ temporal occurrence distributions through the Wasserstein distance to define a proxy distance on the symbol space, used to adapt the Time Warp Edit (TWE) distance to compare symbolic sequences and producing meaningful alignments of sequences.

        \item This distance is leveraged to propose an adaptation of the standard DBA algorithm, yielding rich, compact, and interpretable average representations of "typical" task execution across diagnosis groups. 

        % \item We analyze symbolic sequences via (i) representative descriptive features (e.g., occurrences and durations); (ii) the edit operations-aware distance measure, and (iii) through the computation of average "typical" sequences representing diagnosis-specific executions of the task.

    \end{enumerate}

    \textbf{Contributions:}
    \begin{enumerate}

        %\item We developed a proof-of-concept computational framework to leverage egocentric video recordings for the analysis of complex neurological disorders having subtle behavioral manifestations.% which can ultimately improve our understanding of these conditions and improve patient care.
        
        % \item We propose a novel computational framework for behavioral phenotyping from egocentric video data, based on the computation of average sequences (called barycenters) to represent typical execution of the task per diagnosis groups, holding promises to characterize the behavioral manifestations induced by dysexecutive syndromes, find interpretable and severity-related markers, and improve our understanding of brain lesions-behaviors relationships.

        % \item We illustrate the use of this computational framework for the assessment of  analyze the execution of a complex cooking task in a cohort of 122 participants, including 96 patients with dysexecutive syndromes (37 with RIL and 59 with TBI) and 26 control participants. 

        \item We propose a novel computational framework for behavioral phenotyping from egocentric video data collected during an ecological cooking task, based on descriptive features of the symbolic sequences, the TWE-based distance, and diagnosis-specific barycenters, offering complementary views of task execution variability across severity-stratified diagnosis groups. 

        \item Several variables differed significantly across diagnosis groups, such as the frequency to consult the recipe or the proportion of matching symbols with associated barycenters. This novel approach holds promises to better characterize behavioral manifestations of dysexecutive syndromes, that could ultimately improve clinical assessment and patient care.

        \item Our results show that an approximate solution to the median string problem can be found when inducing a task-dependent metric structure on the symbol space. This allows to solve the global-alignment problem via the efficient DBA algorithms, which were used to compute prototypes of task execution. 

    \end{enumerate}
\end


% filtered

% Question: is the use of the barycenters or the empirical samples the most appropriate choice for defining the symbols temporal occurrences space ? It could be interesting to compute the Wasserstein distance between barycenter-based and empirical-based temporal occurrences distributions across symbols to examine whether they do correspond, to potentially improve our understanding of the average sequence computation method, highlight biases that may appear in the temporal distirbution. It may be that temporal distributions estimated from barycenters corresponds to aligned-versions of symbols occurrences across participants, which may more trusthworty reflect symbols occurrence likelihood, and (i) from which participant's sequences could be compared more reliably, and (ii) yielding a more principled and robust vocabulary distance measure. In future work, personalized analysis could investigate group-differences in symbols temporal occurrences via this vocabulary distance measure. 



% Analysis of symbols is challenging: \cite{vonZiegler2024} Analysis of behavioral flow resolves latent phenotypes

\section{Introduction}
    \label{sec:chapter_6_introduction}

    The computation of symbolic representations of participants egocentric vision during the cooking task offer multiple promising avenues to compare the execution of the task between participants with and without dysexexutive syndromes. We separate the analysis of symbolic sequences across diagnosis groups in two parts. We first compare descriptive variables (e.g., composition, actions durations, number of occurrences) across diagnosis groups, before proposing a novel elastic similarity measure, the Temporal-Wasserstein Time Warp Edit distance (TW-TWE-), used in the computation of average sequences characterizing typical execution of the task.

\section{Descriptive analysis of symbolic representations}
    \label{sec:chapter_6_descriptive_analysis}

    \paragraph{Analyzed cohort.\\}

        From the cohort described in \Cref{sec:chapter_4_analyzed_cohort}, describing 162 administrations of 143 participants, we remove 24 incomplete administrations used to estimate the prototype vocabulary but missing the beginning or end part of the task in the videos, and retain only the first administration for patients with multiple recordings. We thus analyze the symbolic representations of $N=122$ participants, including $n_C=26$ control participants, and $n_P=96$ patients with dysexecutive syndromes (DS). Among patients, $n_{RIL}=37$ have radiation-induced leukoencephalopathy (RIL), and $n_{TBI}=59$ have traumatic brain injury (TBI). We exclude from the analysis $n_{other}=9$ patients with other etiologies to ensure homogeneous patient group. %The demographic and clinical characteristics of the analyzed participants are summarized in \Cref{tab:chapter_6_participants}.


    \paragraph{Analyzed symbolic representations.\\}

    % CODE: Find the proportions
    We use the vocabulary of \texttt{Model-A} to ensure that the distinction between \texttt{task-related} and \texttt{exogeneous} actions corresponds to the definition provided in \Cref{sec:chapter_5_annotation}. Moreover, we address the ambiguity of consolidated prototypes by retaining manually validated symbols only to describe participant's egocentric vision. Specifically, the vocabulary we use combines (i) 54 \textit{valid} prototypes of the $2^{nd}$ consolidation, (ii) 38 prototypes of the $1^{st}$ consolidation for \textit{ambiguous} cases, and (iii) 57 original centroids, in case prototypes from the $1^{st}$ consolidation stage were also classified as \textit{ambiguous} ($N^{1^{st}}_{\text{valid}} = 38/50$; $N^{1^{st}}_{\text{ambiguous}} = 12/50$). The final vocabulary consists in $G=150$ symbols including one "background" segment.

    Given that this final vocabulary is relatively large and contain multiple symbols encoding similar actions (e.g., different presentation of the recipe), we further group them manually into $K=28$ semantically-consistent categories. This grouping allows us to perform a more robust and interpretable analysis of the symbolic sequences, while retaining sufficient granularity to capture the diversity of actions performed during the task. We present this vocabulary and the analyzed symbolic representations in \Cref{fig:chapter_6_vocabulary_SR}.

    \begin{figure}[!htbp]
        \centering
        \includegraphics[width = 1.0\linewidth]{chapters/chapter 6/figures/chapter_6_vocabulary_SR.png}    
        \caption{\textbf{Analyzed symbolic representations of participant's egocentric vision during the cooking task after manual grouping of the prototypes.} (Top) Non-interpolated symbolic representations reflective of real durations, (bottom) interpolated symbolic representations highlighting the temporal consistency and differences between participant's task executions. Each row corresponds to a different participant performing the cooking task, and each color represents the label associated with the embeddings encoding participant's egocentric video. The final vocabulary of $G=150$ prototypes is manually gathered into 28 categories to study participant's behavior during the cooking task across diagnosis groups. Sequences are upsampled to match the longest sequence for visualization purposes.}
        \label{fig:chapter_6_vocabulary_SR}
    \end{figure}

    \subsection{Differences in temporal segmentation of symbolic sequences}
    \label{sec:chapter_6_temporal_Segmentaiton_results}

        \paragraph{Total duration of the task.\\}

        First, we investigate whether the well-known total duration marker in the cooking task differed across clinical groups \cite{Poncet2014, Chevignard2008}. we conducted pairwise Mann–Whitney U tests. Results revealed significant differences between all pairs of groups (p < 0.01), with participants with RIL exhibiting longer durations than both participants with TBI and control, with substantial to moderate effect sizes. This demonstrate that task duration differs systematically across groups, serving as a simple indicator of overall dysexecutive severity. While significant, we seek to compute lower level discriminative variables to capture behavioral differences that manifest during task execution.  % And useful for the analysis! 
        
        \begin{table}[htbp]
            \centering
            \begin{tabular}{l l r r r r r r r}
            \toprule
            \textbf{Group 1} & \textbf{Group 2} & $n_1$ & $n_2$ & Median$_1$ & Median$_2$ & $U$ & $p$-value & $r_{rb}$ \\
            \midrule
                Control & RIL & 26 & 37 & 27.38 & 43.20 & 124.0  & $6.46\times10^{-7}$ & $0.74$ \\
                Control & TBI & 26 & 59 & 27.38 & 33.22 & 434.5  & $1.54\times10^{-3}$ & $0.43$ \\
                RIL     & TBI & 37 & 59 & 43.20 & 33.22 & 1516.5 & $1.40\times10^{-3}$ & $-0.39$ \\
            \bottomrule
            \end{tabular}
            \caption{Mann-Whitney U tests comparing total duration between clinical groups.  $U$ = Mann-Whitney statistic, $p$ = p-value, $r_{rb}$ = rank-biserial effect size. TBI: traumatic brain injury; RIL: radiation-induced leukoencephalopathy.  Significance level was set at $p<0.05$.}
            \label{tab:duration_mannwhitney}

        \end{table}

        \paragraph{Fragmentation of the task and segments duration analysis.\\}

        We investigated whether the fragmentation of participants' activities during the task differed across groups. Fragmentation was quantified as the frequency of segment transition, i.e., the total number of estimated segments (including background segments) divided by the duration of the assessment. The results are summarized in \Cref{tab:fragmentation_results}.

        % \begin{table}[h!]
        %     \centering
        %     \begin{tabular}{l l r r r r r r r}
        %     \hline
        %     Group 1 & Group 2 & $U$ & $p$ & $r_{rb}$ & Median 1 & Median 2 & $n_1$ & $n_2$ \\
        %     \hline
        %     \multicolumn{9}{c}{\textbf{Binary group distinction}} \\
        %     \hline
        %     Control & Patients & 1765.0 & $1.24 \times 10^{-3}$ & -0.414 & 6.0 & 4.9 & 26 & 96 \\
        %     \hline
        %     \multicolumn{9}{c}{\textbf{Subgroup analysis}} \\
        %     \hline
        %     Control & RIL & 760.0  & $1.01 \times 10^{-4}$ & -0.580 & 6.0 & 4.3 & 26 & 37 \\
        %     Control & TBI & 1005.0 & $2.35 \times 10^{-2}$ & -0.310 & 6.0 & 5.2 & 26 & 59 \\
        %     RIL     & TBI & 771.0  & $1.60 \times 10^{-2}$ &  0.294 & 4.3 & 5.2 & 37 & 59 \\
        %     \hline
        %     \end{tabular}
        %     \caption{Pairwise comparisons of fragmentation rate across groups. $U$ = Mann-Whitney statistic, $p$ = p-value, $r_{rb}$ = rank-biserial effect size. Medians indicate the central tendency of fragmentation rate per group. TBI: traumatic brain injury; RIL: radiation-induced leukoencephalopathy.}
        %     \label{tab:fragmentation_results}
        % \end{table}


        \begin{table}[h!]
            \centering
            \begin{tabular}{l l r r r r r r r}
            \hline
            Group 1 & Group 2 & $U$ & $p$ & $r_{rb}$ & Median 1 & Median 2 & $n_1$ & $n_2$ \\
            \hline
            \multicolumn{9}{c}{\textbf{Binary group distinction}} \\
            \hline
            Control & Patients & 1765.0 & $1.24 \times 10^{-3}$ & -0.414 & 6.0 & 4.9 & 26 & 96 \\
            \hline
            \multicolumn{9}{c}{\textbf{Subgroup analysis}} \\
            \hline
            Control & RIL & 760.0  & $1.01 \times 10^{-4}$ & -0.580 & 6.0 & 4.3 & 26 & 37 \\
            Control & TBI & 1005.0 & $2.35 \times 10^{-2}$ & -0.310 & 6.0 & 5.2 & 26 & 59 \\
            RIL     & TBI & 771.0  & $1.60 \times 10^{-2}$ &  0.294 & 4.3 & 5.2 & 37 & 59 \\
            \hline
            \end{tabular}
            \caption{Pairwise comparisons of fragmentation rate across groups. $U$ = Mann-Whitney statistic, $p$ = p-value, $r_{rb}$ = rank-biserial effect size. Medians indicate the central tendency of fragmentation rate per group (transition per minute). TBI: traumatic brain injury; RIL: radiation-induced leukoencephalopathy.  Significance level was set at $p<0.05$.}
            \label{tab:fragmentation_results}
        \end{table}


        % CODE: check that Healthy is not present in the results 

        These results suggest that participants with DS exhibit significantly lower fragmentation rates compared to control participants. Additionally, participants with RIL - which is associated with more severe brain injury - exhibited even lower fragmentation compared to patients with TBI. We further investigate these differences in the durations of non-background segments across diagnosis groups. In \Cref{fig:chapter_6_ecdf}, we visualize the distributions of segments durations across participants using empirical cumulative distribution functions (ECDFs, defined in \Cref{sec:chapter_5_results_human_feedback}).    
            
        \begin{figure}[!htbp]
            \centering
            \includegraphics[width = 1.0\linewidth]{chapters/chapter 6/figures/chapter_6_ecdf.png}    
            \caption{\textbf{Empirical cumulative density function (ECDF) of segments durations.} The ECDF illustrates the distribution of segment durations, with individual participants' ECDFs shown in gray, and group-level ECDFs overlaid for comparison. Shaded areas represent 95\% confidence intervals obtained via bootstrapping. TBI: traumatic brain injury; RIL: radiation-induced leukoencephalopathy.}
            \label{fig:chapter_6_ecdf}
        \end{figure}

        Segment lengths differed significantly across pathology groups (likelihood-ratio test $\chi^2(2)=11.53$, $p=3.31\times10^{-3}$). Relative to the control group, participants with TBI had 8.1\% longer segments, and participants with RIL had 22.9\% longer segments (see statistical methods below). These results corroborate recent findings in the literature, suggesting that low activity fragmentation \citep{10.1001/jamanetworkopen.2021.35168} are associated with mild cognitive impairments \cite{Khosla2025}. %Next, we investigate differences in activity-specific segment durations and number of occurrences across diagnosis groups. %Further investigation is needed to understand the underlying mechanisms driving these differences in fragmentation rates.
        
    \subsection{Comparison of exogeneous vs. task-related actions across groups}
        \label{sec:chapter_6_task_exo_discrimination}

        We hypothesized that patients with dysexecutive syndromes would exhibit a higher proportion of exogenous actions during the cooking task, reflecting difficulties in maintaining focus on task-relevant activities. To test this hypothesis, we classified each prototypical egocentric vision atoms as either \texttt{task-related} or \texttt{exogenous} during the manual annotations described in \Cref{sec:chapter_5_annotation}. We compared the proportions of exogenous segments across groups, calculated as the number of segments classified as exogenous divided by the total number of non-background segments. Pairwise group comparisons were performed using the Mann–Whitney U test, and the rank-biserial correlation was used to estimate effect sizes. The results are summarized in \Cref{tab:task_exo_results}.

        \begin{table}[h!]
            \centering
            \label{tab:task_exo_results}
            \begin{tabular}{l l r r r r r r r}
            \hline
            Group 1 & Group 2 & $U$ & $p$ & $r_{rb}$ & Median 1 & Median 2 & $n_1$ & $n_2$ \\
            \hline
            \multicolumn{9}{c}{\textbf{Binary group distinction}} \\
            \hline
            Control & Patients & 1030.5 & $1.75 \times 10^{-1}$ &  0.174 & 0.083 & 0.11 & 26 & 96 \\
            \hline
            \multicolumn{9}{c}{\textbf{Subgroup analysis}} \\
            \hline
            Control & RIL & 291.0  & $8.15 \times 10^{-3}$ &  0.395 & 0.083 & 0.17 & 26 & 37 \\
            Control & TBI & 739.5  & $7.97 \times 10^{-1}$ &  0.0359 & 0.083 & 0.091 & 26 & 59 \\
            RIL     & TBI & 1499.0 & $2.18 \times 10^{-3}$ & -0.373 & 0.17  & 0.091 & 37 & 59 \\
            \hline
            \end{tabular}
            \caption{Pairwise comparisons in the proportion of exogenous segments across groups. $U$ = Mann-Whitney statistic, $p$ = p-value, $r_{rb}$ = rank-biserial effect size. Medians indicate the central tendency of the proportion of exogenous actions per group. TBI: traumatic brain injury; RIL: radiation-induced leukoencephalopathy. Significance level was set at $p<0.05$.}
        \end{table}

        While no significant difference was observed between control participants and those with acquired brain injuries, the subgroups analysis revealed a more specific pattern: patients with RIL exhibited a significantly higher proportion of exogenous actions (e.g., looking at the window or room, or talking to the examiner) compared to both control participants and patients with TBI. This is likely due to known difficulties in maintaining sustained attention and lack of goal-directed behavior associated with dysexecutive syndromes \cite{Hendry2016sh}. Notably, this marker appeared specific to patients with RIL as no significant difference was found between healthy participants those with TBI.  %We next analyze in detail the composition of participants' symbolic representations for all 28 categories across diagnosis groups.

    \subsection{Comparison of symbolic representations composition across diagnosis groups}

        We examined whether the composition of participants' symbolic representations for all 28 categories differed across diagnosis groups. To this end, we computed three metrics for each of the 28 categories in the vocabulary: (i) the \emph{number of occurrences}, defined as the number of times a category appears in a participant's symbolic representation, (ii) the \emph{mean duration}, defined as the average duration of segments belonging to a category, and (iii) the \emph{proportion of time}, defined as the total duration of segments in a category divided by the total duration of the assessment. We report the subgroups analysis in \Cref{fig:per_label_pathologie_p_values}, and the comparison between control and patients groups in \Cref{apd:chapter_5_control_vs_patients}.

        \begin{figure}[h!]
            \centering
            \includegraphics[width=1.0\textwidth]{chapters/chapter 6/figures/per_label_pathologie_p_values.png}
            \caption{\textbf{Comparison of number of occurrences (top), average duration (middle), and proportion of time (bottom) of each labelled category across diagnosis groups.} Bars indicate mean values per category and group, with participant counts displayed above or below each bar. Error bars represent standard deviation divided by the square root of the sample size (standard error of the mean). Pairwise group differences within each category were tested using the Mann-Whitney U test with Bonferroni-Hochberg correction, and significant differences are shown by brackets connecting the bars with annotated p-values. Statistical significance was set at $p<0.05$. TBI: traumatic brain injury; RIL: radiation-induced leukoencephalopathy.}
            \label{fig:per_label_pathologie_p_values}
        \end{figure}

        These results highlight several notable behavioral differences between participants. For instance, group-level severity of dysexecutive syndromes was negatively associated with the proportion of time spent merging the dough and positively associated with the number of recipe consultations. This corroborates known impairments in episodic or prospective memory leading to a failure to remember the instructions \cite{Hendry2016sh}. Longer durations of oven manipulation appeared specific to the RIL subgroup, which also talked more frequently and for longer periods with the examiner, and required more time to set up the oven compared to patients with TBI. While these analysis demonstrate that symbolic presentations can inform of behavioral differences in task execution across severity-stratified diagnosis groups, further research is needed to understand the mechanisms underlying these findings and explore their clinical implications. %Next, we propose to compare symbolic sequences more holistically using similarity measures and barycenter computation.


    \subsection{Statistical methods}
        Pairwise group comparisons of fragmentation rate, number of occurrences, mean duration, and proportion of time in each category were performed using the non-parametric Mann–Whitney U test, with rank-biserial correlation estimating effect sizes. We corrected for false discovery rate inflation due to multiple comparisons by applying the Bonferroni-Hochberg procedure separately for each metric. To quantify uncertainty around the aggregated ECDFs, we compute $95\%$ confidence intervals via bootstrapping: repeated resampling with replacement of the observed segment durations yields a distribution of ECDF values at each duration, from which the $2.5^{\text{th}}$ and $97.5^{\text{th}}$ percentiles define the shaded confidence bands. Group differences in segment lengths were assessed using a linear mixed-effects model with participant as a random intercept and diagnosis group as a fixed effect. Model comparison was performed via a likelihood-ratio test between the null model (no group effect) and the alternative model (including group) using maximum likelihood estimation. Segment durations were log-transformed to reduce skewness, and fixed-effect estimates were back-transformed to the original scale to report multiplicative changes. This approach accounts for the non-independence of segments within participants and is recommended for hierarchical repeated-measures data \cite{Gelman2006, BAAYEN2008390}. Statistics were calculated in Python V.3.10.16, using SciPy low-level functions V.1.10.1, and Statsmodels V.0.14.4.



\section{Similarity measures and barycenter computation for symbolic representations}
    \label{sec:chapter_6_similarity_measure}
        
    Measuring the distance (or similarity) between symbolic sequences is a primitive operation in many domains, including bioinformatics, natural language processing, and time series analysis \cite{combettes_these}, which can be used to address tasks such as classification, clustering, pattern mining, anomaly detection, and query retrieval \cite{Esling2012}. In the context of egocentric vision symbolic representations, it can provide insights into the commonalities and differences in participants' execution of the task, and quantify inter-individual variability. We consider symbolic sequences defined over a finite alphabet $\mathcal{A}$ of size $G = |\mathcal{A}|$, which consists of the prototype vocabulary indexes. For two symbolic sequences $s_1=(s^1_{1}, ..., s^1_{n}) \in \mathcal{A}^n$ and $s_2=(s^2_{1}, ..., s^2_{m}) \in \mathcal{A}^m$ of length $(m, n)\in \mathbb{N}^2$, we define a distance measure as:

    \begin{equation}
    \begin{aligned}
        D &: \mathcal{A}^m \times \mathcal{A}^n \to \mathbb{R}_+, \\
        & (s_1,s_2) \mapsto D(s_1,s_2),
    \end{aligned}
    \label{eq:distance}
    \end{equation}
    
    quantifying the dissimilarity between the two sequences. The design of such measure should account for temporal distortions between sequences, such as warping (i.e., differences in contraction or dilatation), which may arise in our context from variations in speed or duration, and leading to misalignments and varying sequences lengths. Naive "lock-step" alignment distances such as the Euclidean distance are sensitive to these temporal distortions, and may yield misleading results when comparing sequences that are similar in content but differ in timing \cite{Abanda2018}. 
    
    Elastic distances, such as the popular Dynamic Time Warping (DTW) \cite{Sakoe1971ADP}, have been proposed to address this limitation by allowing for non-linear alignments between sequences during distance computation \cite{combettes_these}. For two D-dimensional time series $\mathbf{x_1} = (x^1_{1}, ..., x^1_{n})^T \in \mathbb{R}^{D \times n}$ and $\mathbf{x_2} = (x^2_{1}, ..., x^2_{m})^T \in \mathbb{R}^{D \times m}$, the DTW is defined as:
    
    $$DTW_p(x_1, x_2) = \min_{\pi \in \Pi(n, m)} \left( \sum_{(i, j) \in \pi} |x^1_{i} - x^2_{ j} |^p \right)^{\tfrac{1}{p}},$$

    where $p \geq 1$ defines the order of the Minkowski distance (with $p=2$ defining the classical Euclidean distance), and $\Pi(n, m)$ is the set of all possible alignment paths between the two time series of length $n$ and $m$. Specifically, for $(n, m) \in \mathbb{N}^2$, an admissible path $P= ((i_1,j_1),...,(i_{L_\pi},j_{L_\pi})) \in (\mathbb{N} \times \mathbb{N} )^{L_\pi}$ of length $L_\pi$ must verify the following conditions:
    \begin{enumerate}
        \item \textbf{Boundary conditions:} $(i_1, j_1) = (1,1)$ and $(i_{L_\pi}, j_{L_\pi}) = (n,m)$.
        \item \textbf{Monotonicity:} $i_{k+1} \ge i_k$ and $j_{k+1} \ge j_k$ for all $k \in [L_\pi - 1]$.
        \item \textbf{Continuity:} $i_{k+1} - i_k \le 1$ and $j_{k+1} - j_k \le 1$ for all $k \in [L_\pi - 1]$.
    \end{enumerate}
    The DTW distance can be computed efficiently using dynamic programming, with a computational complexity of $\mathcal{O}(nm)$ \cite{Sakoe1971ADP}.
    However, DTW does not satisfy the triangle inequality, and thus does not define a proper distance metric \cite{combettes_these}. 
    
    For string sequences, distances typically measure the minimum number of edition (e.g., insertion, deletion, and substitutions for the standard Levenshtein distance) needed to transform one sequence into another \cite{combettes_these}. Other distances on string find Longest Common SubSequences (LCSS) \cite{Hirschberg1977, Needleman1970}, or integrate additional edition operations (e.g., transposition \cite{Damerau1964}, duplication or contraction \cite{Pinhas2013}) to improve robustness to various type of errors or sensitivity to specific patterns of interest. See a recent review of symbolic sequences distance in \cite{combettes_these}. The Time Warp Edit (TWE) distance has been proposed to combine the suitability of edit operations in comparing string sequences, with the benefits of elastic distances to account for temporal misalignments \cite{Marteau2009}. Since the use of edition operations appears a compelling and interpretable choice to compare the execution of a task, for example by modelling addition or absence of a specific behavior, we propose in the following to adapt TWE to our prototype-based symbolic representation approach. 

    \subsection{Time Warp Edit (TWE) distance for symbolic sequences.}

        The Time Warp Edit (TWE) distance searches an optimal alignment path over a cost matrix to account for possible warping in the computation of the distance \cite{Marteau2009}. Given two symbolic representations, potentially of unequal length, $\mathbf{s_1} = (s^1_1, s^1_2, \dots, s^1_n)$ and $\mathbf{s_2} = (s^2_1, s^2_2, \dots, s^2_m)$, TWE computes a cost matrix $C \in \mathbb{R}^{n \times m}$ recursively, which for $1 \leq i \leq n$ and $1 \leq j \leq m$ is defined as follows:
        \begin{align*}
        \text{match} &= C_{i-1,j-1} + d(s^1_i, s^2_j) + d(s^1_{i-1}, s^2_{j-1}) + 2\nu |i-j|, \\
        \text{delete} &= C_{i-1,j} + d(s^1_i, s^1_{i-1}) + \lambda + \nu, \\
        \text{insert} &= C_{i,j-1} + d(s^2_j, s^2_{j-1}) + \lambda + \nu, \\
        D_{i,j} &= \min(\text{match}, \text{delete}, \text{insert}),
        \end{align*}
        where $d(\cdot,\cdot)$ is a pointwise distance function, $\nu\geq 0$ is the stiffness parameter controlling temporal warping, and $\lambda \geq 0$ is the edit penalty applied to both insertion and deletion operations. The final TWE distance between the two series is given by the value $C(n,m)$. The computational complexity of TWE is in $\mathcal{O}(nm)$, and can be computed efficiently using dynamic programming. The stiffness $\nu$ acts as a multiplicative penalty on the temporal misalignment between matched points, with $\nu = 0$ yielding no warping penalty. The edit or gap penalty, $\lambda$, penalizes deletion and insertion of symbols. Importantly, TWE defines a distance metric over the input space, i.e., verifying positivity, separability, symmetry, and the triangle inequality, as long as $\lambda, \nu \geq 0$ and $d(\cdot,\cdot)$ defines a distance metric \cite{Marteau2009}.

        While the original formulation of TWE was developed for real-valued discrete sequences, typically using an $L_p$ distance as the pointwise metric $d(\cdot,\cdot)$, it can be naturally extended to symbolic sequences by defining an appropriate symbol-level distance measure. We propose to construct this distance based on differences in the temporal occurrence patterns of the vocabulary items. Specifically, we measure the distance between symbols using the Wasserstein distance between their normalized temporal distributions, as detailed below. Indeed, we show in the following section that the prescribed nature of the task induces characteristic temporal patterns in the prototypes occurrences, which can be leveraged in the design of this pointwise distance. Notably, we hypothesize that the incorporation of information stemming from video latent representations may introduce arbitrary distortions, as latent space similarities may not accurately order the different symbols compared to the temporal information. 


    
    \subsection{Kernel density estimation (KDE) of prototypes temporal distributions for Wasserstein distance computation}
        \label{sec:chapter_6_kde_temporal_distributions}
        %For example, the initial test setup of \Cref{fig:chapter_4_material} may only be seen at the beginning, reading the recipe typically occur before each new step, and setting up the oven and baking the cake may only be seen at the end of the assessment. Likewise, some visual scene mays occur uniformly without distinction at any time of the assessment, e.g., talking to the examiner, taking a napking, or looking by the window. 


        % Temporal discrimination of prototypes
        %While our initial prototypes are based on spatial clustering, 

        During the cooking task, participants follow the recipe steps in sequential order, starting from the end of the examiner’s explanations and continuing until the cake is placed in the oven. While participants may vary in speed and strategy, the structured nature of the task imposes constraints on the occurrence of egocentric vision atoms. This temporal structure is illustrated in \Cref{fig:chapter_6_vocabulary_SR}, where symbols activations exhibit temporal consistency across participants. To capture these symbol-level temporal occurrences, we compute the empirical probability density function (p.d.f.) of the normalized timestamps (between 0 and 1) associated with each all of their assigned samples. This is performed using a kernel density estimation (KDE) approach, which is a non-parametric method to estimate the probability density function of a random variable, alleviating the injection of priors or pre-specified distribution. For each participant's sequence $s$, let $\mathcal{T}^{(s)}_k \subset [0, 1]^{N^{(s)}_k}$ be the set of $N^{(s)}_k$ normalized timestamps associated with symbol $k \in [G]$. In order to mitigate estimation bias stemming from participants who may contribute disproportionately many timestamps, we adopt a participant-level weighting scheme, described below. This aims to ensure that participants contribute equally to the estimated temporal distribution of each cluster. Specifically, we first compute the contribution of each participant $s$ to cluster $k$, defined as:
        
        \[
        p^{(s)}_k = \frac{|\mathcal{T}^{(s)}_k|}{\sum_{s} \left| \mathcal{T}^{(s)}_k \right|}.
        \]

        To correct for potential high dominance of the contribution of participants, we define a normalized weight for each participant $s$ and cluster $k$ as:
        \[
            w^{(s)}_k = \sqrt{ \frac{1}{p^{(s)}_k}}.
        \]
        The full set of weighted timestamps for cluster $k$ is then:
        \[
        \mathcal{S}_k = \bigcup_s \left\{ (t, w^{(s)}_k) \,:\, t \in \mathcal{T}^{(s)}_k \right\}.
        \]

        Finally, we estimate the temporal distribution $\tau_k(t)$ of element $k$ using a weighted kernel density estimator:
        \[
        \hat{\tau}_k(t) = \frac{1}{Z_k h} \sum_{(t_i, w_i) \in \mathcal{S}_k} w_i \cdot \kappa\left( \frac{t - t_i}{h} \right),
        \]
        where $\kappa$ is a positive semi-definite (p.s.d.) kernel, $h$ is the kernel bandwidth, and $Z_k$ is a normalization constant ensuring that $\int_0^1 \hat{f}_k(t) \, dt = 1$. This KDE captures the normalized temporal concentration of each cluster, while correcting for sampling imbalance across participants. We present the estimated temporal distributions of the analyzed vocabulary in \Cref{fig:chapter_6_temporal_distributions_examples}, which suggests that the egocentric vision prototypes effectively captures the temporal structure of the task. In the following, the hat notation is omitted for simplicity, and we denote the estimated temporal distributions as $\tau^{\mathcal{V}} = (\tau_k)_{k \in [G]}$.

        \begin{figure}[!htbp]
            \centering
            \includegraphics[width = 0.9\linewidth]{chapters/chapter 6/figures/chapter_6_temporal_distances_distributions.png}    
            \caption{\textbf{Empirical probability density functions of the vocabulary elements' temporal occurrences (normalized).} A variety of temporal patterns can be observed, including uniform (e.g., background, cleaning hands or talking to the examiner), unimodal (e.g., adding ingredients, opening the fridge, etc.), or multi-modal (e.g., reading the recipe or using the oven). Symbols are sorted by average temporal occurrence.}
            \label{fig:chapter_6_temporal_distributions_examples}
        \end{figure}


        \paragraph{Comparing temporal distributions via Wasserstein (Earth Mover's) distance.\\}            

        We compute the distance between the temporal occurrence patterns of vocabulary items via the Wasserstein distance (also known as the Earth Mover's Distance) \cite{peyreCuturi, Kantorovich1960}. For two one dimensional temporal distributions $\tau_k$ and $\tau_l$, with respective cumulative distribution functions (CDFs) \(F_{\tau_k}\) and \(F_{\tau_l}\), the Wasserstein distance is defined as:
        \[
        W_p(\tau_k, \tau_l) = \int_{0}^{1} \left| F_{\tau_k}(t) - F_{\tau_l}(t) \right|^p \, dt.
        \]

        %, which captures both the difference in temporal location (i.e., when distributions are shifted in time) and the difference in shape (e.g., spread or multi-modality). 
        This distance has several desirable properties for our application: (i) it is a true distance metric, satisfying positivity, symmetry, and the triangle inequality; (ii) it is sensitive to both temporal shifts and shape differences between distributions; and (iii) it can be computed efficiently in one dimension.
        In our experiments, we use the Wasserstein-1 distance ($p=1$) computed from the sampled timestamps and associated weights. We report the empirical Wasserstein distance matrix in \Cref{fig:wasserstein_distance_computation}-a, with vocabulary elements ordered by average temporal occurrence.
        
        \begin{figure}[!htbp]
            \centering
            \includegraphics[width = 0.9\linewidth]{chapters/chapter 6/figures/wasserstein_distance_computation.png}    
            \caption{\textbf{(a) Wasserstein distance between vocabulary elements' temporal distributions and (b) transformed Wasserstein distance used in the TWE distance computation.} (a) The Wasserstein distance captures both temporal shifts and shape differences between distributions, with smaller distances along the diagonal indicating closer temporal occurrence patterns. (b) The final pointwise distance matrix incorporates an offset $\delta=0.3$ for off-diagonal elements to favor matching during TWE distance computation, scale distances between 0 and 1, and majorize the cost associated with the background label.}
            \label{fig:wasserstein_distance_computation}
        \end{figure}

    \subsection{Definition of the temporal Wasserstein TWE distance (TW-TWE) for symbolic sequences of task execution}

    
        To define the pointwise distance measure between symbols $d(\cdot,\cdot)$ used in the TWE distance, we slightly modify the Wasserstein distance presented above. Specifically, we apply in this order: (i) the addition of an offset $\delta$ to off-diagonal elements, to favor alignment paths matching identical symbols by increasing the cost of matching different symbols, (ii) min-max scaling, to facilitate comparison and interpretation,  and (iii) major the cost associated with the background label, by using the symbol-level maximum distance, to prevent excessive alignment with the background category. For two vocabulary elements $(k,l) \in \mathcal{A}^2$, we note $d_{\text{TW-TWE}}(k,l)$ the distances defined by the slightly modified Wasserstein distance described above (including the $\delta$ offset, scaling, and background cost adjustment).

        % This defines a pointwise distance matrix $D_{\text{TW-TWE}} \in \mathbb{R}^{G \times G}$, with elements:
        % \[
        % d_{\text{TW-TWE}}(k,l) =
        % \begin{cases}
        % 0 & \text{if } k = l, \\
        % \frac{W_1(\tau_k, \tau_l) + \delta}{\max_{i,j \in [G]} (W_1(\tau_i, \tau_j) + \delta)} & \text{if } k \neq l \text{ and } k,l \neq \text{background}, \\
        % \frac{W_1(\tau_k, \tau_l) + \delta}{\max_{i \in [G], j \neq \text{background}} (W_1(\tau_i, \tau_j) + \delta)} & \text{if } k = \text{background} \text{ or } l = \text{background}.
        % \end{cases}
        % \]
        Our distance measure between two symbolic representations $\mathbf{s_1}$ and  $\mathbf{s_2}$, potentially of unequal length, is thus defined by using the $d_{\text{TW-TWE}}(\cdot,\cdot)$ as the pointwise distance function in the original TWE distance formulation. We call this distance the TW-TWE distance for clarity, standing for temporal Wasserstein-TWE distance. It is defined as $C_{\text{TW-TWE}}(n,m)$, where the cost matrix $C_{\text{TW-TWE}} \in \mathbb{R}^{n \times m}$ is computed recursively as follows:

        For $1 \leq i \leq n$ and $1 \leq j \leq m$:
        \begin{align*}
        \text{match} &= C^{\text{TW-TWE}}_{i-1,j-1} + d_{\text{TW-TWE}}(s^1_i,s^2_j) + d_{\text{TW-TWE}}(s^1_{i-1}, s^2_{j-1}) + 2\nu |i-j|, \\
        \text{delete} &= C^{\text{TW-TWE}}_{i-1,j} + d_{\text{TW-TWE}}(s^1_i, s^1_{i-1}) + \lambda + \nu, \\
        \text{insert} &= C^{\text{TW-TWE}}_{i,j-1} + d_{\text{TW-TWE}}(s^2_j, s^2_{j-1}) + \lambda + \nu, \\
        D_{i,j} &= \min(\text{match}, \text{delete}, \text{insert}).
        \end{align*}
        
        We detail the hyperparameters selection procedure (for $\delta$, $nu$, and $\lambda$) in \Cref{sec:chapter_6_wasserstein_distance_matrices}, and visualize the distance matrix used in our experiments in \Cref{fig:wasserstein_distance_computation}-b, with $\delta=0.3$.

        % \paragraph{Properties of the TW-TWE distance.\\}
        % The TW-TWE distance inherits the desirable properties of the original TWE distance, including being a true distance metric, as long as the pointwise distance $d_{\text{TW-TWE}}(\cdot,\cdot)$ satisfies the properties of a distance metric. The incorporation of the Wasserstein distance between temporal distributions as the pointwise distance measure allows the TW-TWE distance to capture both temporal misalignments and differences in the temporal occurrence patterns of symbols. The hyperparameters $\nu$, $\lambda$, and $\delta$ provide flexibility in controlling the sensitivity of the distance to temporal warping, edit operations, and symbol matching, respectively. This makes the TW-TWE distance a versatile tool for comparing symbolic sequences in the context of egocentric vision task execution.


        \paragraph{Definition of the proportion of matching symbols.\\}
            For a pair of sequences $s_1 \in \mathcal{A}^n$ and $s_2 \in \mathcal{A}^m$ of length $n$ and $m$, and their optimal alignment path $\pi^* = ((i_1,j_1),...,(i_{L_{\pi^*}},j_{L_{\pi^*}}))$ obtained from the TW-TWE distance computation, we define the proportion of matching symbols as:
            \[
            p_{\text{match}} = \tfrac{1}{L_{\pi^*}} \sum_{(i,j) \in \pi^*} \mathbbm{1}(s^1_i = s^2_j),
            \]
            where $\mathbbm{1}(\cdot)$ is the indicator function. This measure quantifies the similarity in symbolic content after aligning participant's sequences, regardless of temporal distortions and sequence length differences. %This proportion of matches is used in \Cref{sec:chapter_6_wasserstein_distance_matrices} to evaluate the quality of alignments obtained with different TW-TWE hyperparameter settings.

        \paragraph{Implementation details.\\}
            Wasserstein distances between temporal distributions are estimated using the \texttt{ot.emd2\_1d} function of the \texttt{POT} Python library, with a squared euclidean metric $'metric=sqeuclidean'$. We use the 'KDEUnivariate' class of the \texttt{statsmodels} Python library to estimate the temporal distributions, with a Gaussian kernel and a bandwidth $h$ set via Silverman's rule of thumb \cite{silverman1986density}. Adaptation of the TWE distance is based on the \texttt{aeon} Python library \cite{aeon_library}. The code will be made available on github in the wake of the thesis. 
                    
            
    \subsection{Pairwise TW-TWE distance matrix between participant's symbolic sequences}
    \label{sec:chapter_6_wasserstein_distance_matrices}

        The TW-TWE distance involves three hyperparameters: the stiffness $\nu$, the edit penalty $\lambda$, and the offset $\delta$ used in the pointwise distance matrix. To select appropriate values for these hyperparameters, we employed a grid search approach, evaluating the quality of the resulting alignments between symbolic sequences.  Specifically, we aimed to identify hyperparameter values that yield alignments with a high proportion of symbols matches, indicating that the sequences are well-aligned in terms of their symbolic content. We report matches proportion over the following range of values: $\nu \in \{10^{-4}, 10^{-5}\}$, $\lambda \in \{0.0, 10^{-3}, 10^{-2}, 10^{-1}\}$, and $\delta \in \{0, 0.1, 0.2, 0.3, 0.4, 0.5\}$. For each combination of parameters, we computed the pairwise TWE distances between 2000 randomly samples pairs of sequences ($\approx 10\%$ of all pairs), and record the edit operations associated with the optimal alignment paths. We selected the combinations of hyperparameters that yielded the largest proportion of matches when comparing the sequences. We report the grid search performances in \Cref{fig:chapter_6_hyperparameters_searches}, and used $\nu=10^{-5}$, $\lambda=0.1$, and $\delta=0.3$ for TW-TWE distance computation in our experiments.

        \begin{figure}[!htbp]
            \centering
            \includegraphics[width = 1.0\linewidth]{chapters/chapter 6/figures/chapter_6_hyperparameters_searches.png}    
            \caption{\textbf{Boxplot of the percentage of matches between sampled pairs of participant's symbolic sequences using the Time Warp Edit (TWE) distance with grid search hyperparamers.} Each boxplot summarizes the distribution of matches percentage over 2000 randomly sampled pairs of sequences for each combination of hyperparameters. The selected hyperparameters, yielding the highest average percentage of matches, that is, $\nu=10^{-5}$, $\lambda=0.1$, and $\delta=0.3$.}
            \label{fig:chapter_6_hyperparameters_searches}
        \end{figure}

        We visualize in \Cref{fig:chapter_6_examples_matching} illustrative examples of optimal alignments, TW-TWE distance, and percentage of matching symbols between participant's symbolic sequences using the selected hyperparameters. The TW-TWE distance appears to capture the commonalities between sequences, as indicated by the high proportion of matches present in the optimal alignments between sequences, which suggests it could be a useful tool to compare participants' execution of the task. %We next compute the pairwise TWE distance matrix between all participants' symbolic representations, and investigate its structure to identify potential clusters of similar participants.

        \begin{figure}[!htbp]
            \centering
            \includegraphics[width = 1.0\linewidth]{chapters/chapter 6/figures/chapter_6_examples_matching.png}    
            \caption{\textbf{illustrative examples of optimal alignments between participant's symbolic sequences using the TW-TWE distance.} The TW-TWE distance is computed between participant's symbolic representations using optimal hyperparameters $\hat{\nu}=10^{-5}$, $\hat{\lambda}=0.1$, and $\hat{\delta}=0.3$. Pairs of symbols along the optimal TWE alignment are colored to highlight matches between symbols (green), mismatch (i.e., substitution; red), and matches between background symbols (light green).}
            \label{fig:chapter_6_examples_matching}
        \end{figure}

        \paragraph{Comparison of task execution using the TW-TWE distance over participants symbolic sequences.\\}
        
        We compute and visualize the pairwise TW-TWE distance between participant's symbolic representations in \Cref{fig:chapter_6_distance_matrix_pairwise}-a, sorted by diagnosis groups to facilitate interpretation. In order to alleviate the distance dependence to sequence length, which was proven to differ significantly between diagnosis groups in \Cref{sec:chapter_6_temporal_Segmentaiton_results}, we interpolate sequences to the median sequence length of the analyzed cohort, that is, 5182 symbols ($\approx 27.4$ minutes). This is argued to yield no significant differences in performances in the literature \cite{Ratanamahatana2005}, while enabling fair comparisons of distance across participants. We also perform hierarchical clustering using Ward's method on the distance matrix to identify potential clusters of similar participants, and visualize the resulting dendrogram alongside the ordered distance matrix in \Cref{fig:chapter_6_distance_matrix_pairwise}-b.

        \begin{figure}[!htbp]
            \centering
            \includegraphics[width = 1.0\linewidth]{chapters/chapter 6/figures/chapter_6_distance_matrix_pairwise.png}    
            \caption{\textbf{Pairwise TWE distance matrix sorted per diagnosis groups (a) and ordered per hierarchy of distances with associated dendogram and diagnosis groups (b).} }
            \label{fig:chapter_6_distance_matrix_pairwise}
        \end{figure}

        As can be seen, the distance between sequences alone does not appear to clearly separate diagnosis groups, as indicated by the intertwining of colors in the sorted distance matrix. Given the TW-TWE were found to align participant's sequences effectively, this suggests that the level of inter-individual variability in task execution may prevail over group-level differences, making it challenging to differentiate between diagnosis groups based solely on warping aware distances. Further refinement of a similarity measure between symbolic sequences, for example invariant to symbols ordering, may be investigated in future work to better capture group-level differences. 
        
        
        %further refinement of the symbolization procedure and of the distance computation may help to discriminate between diagnosis groups more effectively. %Next, we propose to leverage these optimal alignments for computing average symbolic representation of each diagnosis groups.
        
        
        %further refinement of the symbolization procedure and of the distance computation may help to discriminate between diagnosis groups more effectively. %Next, we propose to leverage these optimal alignments for computing average symbolic representation of each diagnosis groups.

\section{Computation of average symbolic representations (barycenters) for each diagnosis group}
 \label{sec:chapter_6_barycenter_computation}   

    Gold standard assessment of participant’s execution of the cooking task consists in identifying errors performed during the task, such as the ones described in \Cref{apd:chapter_4_error_type_description}, using manual coding of participant's behavior. Their complexity and subtlety make it challenging to identify these errors automatically from symbolic representations. Instead, we hypothesized that the computation of \textit{average} symbolic sequences could help to identify both the common and specific aspects of participant's behaviors during the cooking task, using a quantified and interpretable quantity. 

    In our context, we aim to estimate average sequences that could represent typical egocentric vision trajectories during the cooking task, that is, temporally ordered to be consistent with the order of recipe steps, and with symbols durations representative of the empirical average duration of each steps. We believe that such average sequences, not necessarily uniquely defined, could provide useful reference points to be compared with participant's sequences, holding promises to reveal characteristic differences in task execution associated with different diagnosis groups. 
    We thus ask the questions of (i) whether the computation of average egocentric vision sequences during the cooking task is feasible; (ii) if a comparison of barycenters of each diagnosis group can effectively reveal characteristic differences in task execution, and (iii) whether the comparison of participant's symbolic representation to barycenters can differentiate between severity-stratified diagnosis groups.
    

    %These can define prototypical representation of task execution for each diagnosis group, potentially revealing insights into the characteristic execution patterns associated with each pathology. These may ultimately be used to detect automatically and robustly patterns of errors that are common in executive functioning impairments. 

    Our egocentric vision vocabulary primarily defines categorical symbols, without inherent ordering or metric structure, thus defining our symbolic sequences as string, where finding an average over an edit distance typically refers to the "median string problem" \cite{KOHONEN1985309}. This is a well-studied problem in bioinformatics, which is known to be NP-hard when the alphabet is unbounded \cite{DELAHIGUERA200039}, and W[1]-hard with respect to string length even for finite alphabets \cite{NICOLAS2005390} (see \cite{fine_grained_complexity} for a recent fine-grained complexity analysis). As a result, computing the exact median string is generally intractable, motivating the use of heuristic procedures that iteratively refine an initial candidate through local edit operations until a satisfactory approximation is reached \cite{Mirabal2019PivotSF}.
    
    However, given that we induce a proxy ordering for egocentric vision symbols, based on their typical temporal occurrences, we study the computation of average symbolic representation leveraging the methodology of time series barycenters. Specifically, we use the state-of-the-art Dynamic Time Warping Barycenter Averaging (DBA) algorithm \cite{PETITJEAN2011678}, which we adapt to the symbolic sequences domain, notably via the TW-TWE distance.



    % We discuss this limitation further below when discussing locally optimal solution estimates. See \cite{Brning2025} for additional convergence and optimality results of DBA.

    % We propose to define the average sequences of a set, or barycenter, as the Frechet mean under TW-TWE distance, that is, the sequence minimizing the sum of TW-TWE distances to all sequences in the set.
    
    % Next, we detail our methodology to compute average symbolic representation using our adapted TWE distance and the standard Dynamic Time Warping Barycenter Averaging (DBA) algorithm \cite{PETITJEAN2011678}.

    % This shows that a plausible solution of the median string problem can be found when inducing a metric space in the space of symbols, such that  as a metric space
    
    % We propose to compute this reference symbolic representation as the barycenter of control participants’ symbolic sequences. Likewise, the barycenter of each diagnosis group could provide a prototypical representation of the task execution for each group, providing insights into the characteristic execution patterns associated with each pathology.
    

    \subsection{Adaptation of the Dynamic Time Warping Barycenter Averaging (DBA) algorithm for symbolic sequences}
    \label{sec:chapter_6_adaptation_DBA}

    Barycenters can act as reference points to compare between participant-level sequences with set-level prototypes, with which differences in alignment and composition may provide a robust quantification of task execution differences. For clustering algorithms such as $k$-means, which rely on the computation of centroids to represent clusters, finding a meaningful average is crucial, and is seen as an higher stake objective compared to the definition of the similarity measure itself \cite{fawaz2023shapedba}. The Dynamic Time Warping Barycenter Averaging (DBA) algorithm is a standard heuristic used for computing the average of time series data \cite{10.1109ICDM.2014.27}, overcoming the limitation of iterating over pairwise sequences by solving a global alignment problem at each iteration \cite{PETITJEAN2011678}. The estimated mean sequence is defined as the minimizer of the sum of DTW with all training sequences. DBA iteratively refines an initial sequence by (i) finding for each of its timestamps their best aligned timestamps with all the sequences in the dataset, then (ii) updating the value at each timestamp as an \emph{average} of their aligned values. The initial barycenter is typically chosen as a random sequence from the dataset, and the algorithm iterates until convergence, defined via the sum of distance between the barycenter and the dataset sequences, or after a maximum number of iterations. The elastic distance used to compute the optimal alignments varies, using for example the DTW as in the original paper, or more refined variants improving performances. Examples include the differentiable softDTW \cite{10.5555/3305381.3305474}, or the more robust neighborhood-aware ShapeDTW \cite{zhao2018shapedtw}, whose adaptation to DBA (refers to as ShapeDBA;  \cite{fawaz2023shapedba}) was shown to improve performances across classification and clustering tasks \cite{fawaz2023shapedba}. 

    While the \emph{average} at each timestamp is typically computed as the mean of the aligned values, it needs adaptation to derive an admissible barycenter for string sequences. For strings, we propose to compute the updated value at each timestamps as the dominant mode of their aligned symbols, which guaranties the admissibility of the estimated barycenter at each iteration. This introduces a slight memory cost overhead to store an accumulation matrix of shape $G\times L$ at each iteration, which appears negligible in practice (e.g. in comparison to the input sequences memory complexity in $\mathcal{O}(N \times L)$). We present this adaptation of the DBA algorithm in \Cref{alg:chapter_6_dba_algorithm}.

    

    % \begin{algorithm}
    %     \caption{Dynamic Time Warping Barycenter Averaging (DBA) for symbolic sequences under the Wasserstein-based TWE distance.}
    %     \label{alg:chapter_6_dba_algorithm}
    %     \begin{algorithmic}[1]
    %         \Require $S = (\mathbf{s}^{(1)},\dots,\mathbf{s}^{(N)})$ a set of $\mathcal{V}$-valued symbolic sequences, $W = (w^{(1)},\dots,w^{(N)})$ the sequence weights, $D_{\mathcal{V}} \in \mathbb{R}^{|\mathcal{V}| \times |\mathcal{V}|}$ the Wasserstein distance matrix between symbols' normalized temporal distribution, $\mathbf{s}_{ini}$ (optional) the initial sequence, $L$ the length of the averaging sequence, $n_{it}$ the number of iterations. 
    %         \If{$\mathbf{s}_{ini} \text{ is given}$}
    %         \State $\mathbf{s}^{(0)} \gets \mathbf{s}_{ini}$
    %         \Else
    %         \State $\textit{Initialize } \mathbf{s}^{(0)} \in S$
    %         \EndIf
    %         \State $\textit{Initialize best solution } s^* \gets s^{(0)}$
    %         \State $\textit{Initialize cost } c^* \gets + \infty$

    %         \For{$iter=1,\dots,n_{it}$}
                
    %             \State $\textit{Initialize accumulator } A \in \mathbb{0}^{G \times L}$
    %             \State $\textit{Initialize cost} c_{iter}=0$
    %             \For{$i=1,\dots,N$}
    %                 \State $\pi^{(iter)} \gets \textit{Optimal TW-TWE alignment path between $\mathbf{s}^{(iter-1)}$ and $s^{(i)}$}$
    %                 \State $d_{TW-TWE}^{(i)} \gets \textit{TW-TWE distance between $\mathbf{s}^{(iter-1)}$ and $s^{(i)}$}$
    %                 \State $\textit{Update accumulator:}$
    %                 \For{$(i, j) \in \pi^{(iter)}$}
    %                     \State $A_{s_j, i} \gets A_{s_j, i} + w^{(i)}$
    %                     \State $c_{iter} \gets c_{iter} + d_{TW-TWE}^{(i)}* w^{(i)}$
    %                 \EndFor 
    %             \EndFor
    %             \State $\textit{Update average symbolic representation:}$
    %             \For{$k=1,\dots,L$}
    %                 \State $\mathbf{s}^{(iter)}_k \gets \argmax_{g \in [G]} A_{g,k}$
    %             \EndFor
    %             % Check best cost solution
    %             \If{$ c_{iter} < c^*$}
    %                 \State $\textit{Update best solution: } \mathbf{s}^* \gets \mathbf{s}^{(iter)}$
    %                 \State $\textit{Update best cost: } c^* \gets c_{iter}$
    %             \EndIf
        
    %         \EndFor
    %         \State \Return $\mathbf{s}^*$
    %     \end{algorithmic}
    % \end{algorithm}


    % \begin{algorithm}
    % \caption{DBA with TW-TWE adapted to symbolic sequences}
    % \label{alg:chapter_6_dba_algorithm}
    % \begin{algorithmic}[1]
    % \Require $S = (\mathbf{s}^{(1)},\dots,\mathbf{s}^{(N)})$ a set of $\mathcal{V}$-valued sequences of length $L$, 
    %         $W = (w^{(1)},\dots,w^{(N)})$ sequence weights, 
    %         $D_{\mathcal{V}} \in \mathbb{R}^{|\mathcal{V}| \times |\mathcal{V}|}$ Wasserstein symbol distance, 
    %         optional initial sequence $\mathbf{s}_{ini}$, averaging length $L$, iterations $n_{it}$
    % \If{$\mathbf{s}_{ini} \text{ is given}$}
    %     \State $\mathbf{s}^{(0)} \gets \mathbf{s}_{ini}$
    % \Else
    %     \State Initialize $\mathbf{s}^{(0)} \in S$
    % \EndIf
    % \State $\mathbf{s}^* \gets \mathbf{s}^{(0)}$, $c^* \gets +\infty$
    % \For{$iter = 1,\dots,n_{it}$}
    %     \State $A \gets \mathbf{0}_{G \times L}$ \Comment{accumulator}
    %     \State $c_{iter} \gets 0$
    %     \For{$i = 1,\dots,N$}
    %         \State Compute optimal TW-TWE alignment $\pi^{(iter)}$ between $\mathbf{s}^{(iter-1)}$ and $\mathbf{s}^{(i)}$
    %         \State Compute $d_{TW-TWE}^{(i)}$ distance
    %         \For{$(p,q) \in \pi^{(iter)}$}
    %             \State $A_{s_q, p} \gets A_{s_q, p} + w^{(i)}$
    %             \State $c_{iter} \gets c_{iter} + w^{(i)} \, d_{TW-TWE}^{(i)}$
    %         \EndFor
    %     \EndFor
    %     \For{$k=1,\dots,L$}
    %         \State $\mathbf{s}^{(iter)}_k \gets \arg\max_{g \in [G]} A_{g,k}$
    %     \EndFor
    %     \If{$c_{iter} < c^*$}
    %         \State $\mathbf{s}^* \gets \mathbf{s}^{(iter)}$, $c^* \gets c_{iter}$
    %     \EndIf
    % \EndFor
    % \State \Return $\mathbf{s}^*$
    % \end{algorithmic}
    % \end{algorithm}


    \begin{algorithm}[ht]
    \caption{DBA for symbolic sequences under the Wasserstein-based TWE distance}
    \label{alg:chapter_6_dba_algorithm}
    \KwIn{$S = (\mathbf{s}^{(1)},\dots,\mathbf{s}^{(N)})$ a set of $\mathcal{V}$-valued symbolic sequences, $W = (w^{(1)},\dots,w^{(N)})$ the sequence weights, $D_{\mathcal{V}} \in \mathbb{R}^{|\mathcal{V}| \times |\mathcal{V}|}$ the Wasserstein distance matrix between symbols' normalized temporal distribution, $\mathbf{s}_{ini}$ (optional) the initial sequence, $L$ the length of the averaging sequence, $n_{it}$ the number of iterations.}
    
    \KwOut{$\mathbf{s}^*$, the DBA average sequence}

    \If{initial sequence $\mathbf{s}_{ini}$ is given}{
        $\mathbf{s}^{(0)} \gets \mathbf{s}_{ini}$\;
    }
    \Else{
        Initialize $\mathbf{s}^{(0)} \in S$\;
    }

    $\mathbf{s}^* \gets \mathbf{s}^{(0)}$, $c^* \gets +\infty$\;

    \For{$iter = 1$ \KwTo $n_{it}$}{
        $A \gets \mathbf{0}_{G \times L}$ \tcp*{accumulator}
        $c_{iter} \gets 0$\;
        \For{$i = 1$ \KwTo $N$}{
            $\pi^{(i)} \gets \textit{Optimal TW-TWE alignment path between $\mathbf{s}^{(iter-1)}$ and $s^{(i)}$}$\;
            $d_{TW-TWE}^{(i)} \gets \textit{TW-TWE distance between $\mathbf{s}^{(iter-1)}$ and $s^{(i)}$}$\;
            \ForEach{$(p,q) \in \pi^{(i)}$}{
                $A_{s_q, p} \gets A_{s_q, p} + w^{(i)}$ \tcp*{update accumulator}
                $c_{iter} \gets c_{iter} + w^{(i)} \, d_{TW-TWE}^{(i)}$\;
            }
        }
        \For{$k = 1$ \KwTo $L$}{
            $\mathbf{s}^{(iter)}_k \gets \arg\max_{g \in [G]} A_{g,k}$ \tcp*{Update average symbolic representation}
        }
        \If{$c_{iter} < c^*$}{
            $\mathbf{s}^* \gets \mathbf{s}^{(iter)}$, $c^* \gets c_{iter}$\;
        }
    }
    \Return{$\mathbf{s}^*$}
    \end{algorithm}

    \paragraph{Implementation details.\\}
        Adaptation of the DBA algorithm is based on the \texttt{aeon} Python library \cite{aeon_library}. The code will be made available on github in the wake of the thesis. 


    \subsection{Experimental setting}

        We compute a barycenter symbolic representation for each diagnosis group with the following settings. We split group-level sequences into $50\%$ training sequences from which computing an average sequence, and $50\%$ testing sequences, repeating the experiment 10 times while ensuring variable splits definition. For each split, we estimate 3 barycenters using different initial sequence randomly sampled. This yields a total of 30 estimated barycenters per group. Additionally, (i) we set the length of the barycenter to the median sequence length, as in \Cref{sec:chapter_6_wasserstein_distance_matrices}, (ii) we set the number of iterations to $n_{it}=50$, which was empirically found to be sufficient for convergence, and (iii) we set uniform weights for all sequences within the training sequences of the group. We interpolate each sequences to the median sequence length to alleviate the dependence of the distance to sequence lengths. Additionally, we tested including or excluding the background symbol when updating the barycenter at each iteration, finding that excluding it can yield more informative barycenters and classification performances (see more information below). Finally, we compute barycenters estimated (i) using the original TWE distance, to further assess the benefit of our TW-TWE distance adaptation, and (ii) using a majority voting prediction, applied to each timestamps of the interpolated symbolic representations for each group. This allows to assess the benefit of using an elastic distance to compute barycenters, compared to a non-elastic approach.
        
        For evaluating the quality of the estimated barycenters, we compute their TW-TWE distance with all sequences in the dataset, and record the proportion of matching symbols associated with optimal alignment paths. We compare TW-TWE distances and proportions of matches across diagnosis groups, to assess whether the estimated barycenters capture characteristic differences in task execution between groups. Statistical comparisons are performed using the Mann-Whitney U tests \cite{mwu} between distances and proportions of matches across groups and between training and testing sets, following the illustration of \Cref{fig:chapter_6_comparison_barycenter}. Effect sizes are computed via the the rank biserial correlation \cite{biserial}. We also visualize the estimated barycenters for each diagnosis group to qualitatively assess their differences. Finally, we investigate the potential of using the distances (respectively, the proportions of matches) to the barycenters in differentiating across diagnosis groups by assigning each sequence its closest barycenter in terms of distance (respectively, highest proportion of matches). We evaluate diagnostic accuracy for pairwise classification of diagnosis groups by computing the area under the receiver operating characteristic curve (AUC-ROC) for each pairwise comparison, with 95\% confidence intervals (CIs) estimated using the Hanley McNeil method \cite{Hanley1982}.

        % Code: caption
        \begin{figure}[!htbp]
            \centering
            \includegraphics[width = 1.0\linewidth]{chapters/chapter 6/figures/chapter_6_comparison_barycenter.png}    
            \caption{\textbf{Overview of the comparisons of barycenters to participant's symbolic sequences.} For each diagnosis group, we compute a barycenter using the DBA algorithm on $50\%$ of the group's sequences (training set), and evaluate its TW-TWE distance and proportion of matching symbols with participants sequences following the four types of comparisons described in the figure.} 
            \label{fig:chapter_6_comparison_barycenter}
        \end{figure}


    \subsection{Results}

        \subsubsection{Comparison of barycenters across diagnosis groups}

        % Providing that estimated barycenters accurately summarize group-level symbolic sequences, they could provide a valuable tool to characterize typical differences in task execution across groups via barycenters estimation, complementary to the descriptive empirical analysis of \Cref{sec:chapter_6_descriptive_analysis}.
        %CODE: refine the theory here
        % We analyze the performance of the proposed DBA algorithm in estimating representative symbolic sequences for each diagnosis group.
        We first analyze the composition of the barycenters across diagnosis groups, in particular compared to the descriptive composition analysis of group-level symbolic sequences in \Cref{sec:chapter_6_descriptive_analysis}, to assess whether the estimated barycenters capture characteristic differences in task execution between groups. For conciseness, we report the analysis on background-excluded barycenters only, as they were found to degrade pairwise classification results. We visualize the estimated barycenters for each diagnosis group in \Cref{fig:chapter_6_non_bg}, and report the estimated background-included barycenters in \Cref{apd:chapter_6_barycenters_background}. Additionally, we report the estimated barycenters using the original TWE distance and the the majority voting baseline method in \Cref{apd:chapter_6_barycenters_median_twe}, which both yield degenerated barycenters due to the arbitrary pairwise distance between symbols, and lack of elastic alignment between sequences, respectively.

    
        \begin{figure}[!htbp]
            \centering
            \includegraphics[width = 1.0\linewidth]{chapters/chapter 6/figures/chapter_6_non_bg.png}    
            \caption{\textbf{Estimated barycenters for each diagnosis group.} Each barycenter is computed using the DBA algorithm under the TW-TWE distance, with the background symbol excluded when updating the barycenter at each iteration.}
            \label{fig:chapter_6_non_bg}
        \end{figure}

        %We noticed that although DBA converges after only a few iterations (typically less than 20)
        We observe qualitatively that the estimated barycenters appear as reasonable average sequences of participants' empirical sequences (e.g., compared to participants' sequences in \Cref{fig:chapter_6_vocabulary_SR}). Moreover, the estimated barycenters vary across both split definitions (10 iterations per diagnosis group) and initial conditions (3 per split definition), highlighting both within-group variability and the sensitivity of the estimated barycenter to initial conditions. Barycenters appear concordant with the results of \Cref{sec:chapter_6_descriptive_analysis}. For example, patient group barycenters exhibit more frequent recipe consultations (purple color), with a longer proportion of total duration dedicated to the initial reading barycenters. Moreover, RIL barycenters exhibit longer average duration for incorporating the white eggs (grey color) than other groups, and control group barycenters exhibit larger proportions of time merging the dough (beige color). However, differences between previous analysis can also be observed, for example the TBI group barycenter manifests a clear pattern of larger relative duration for the "opening the fridge" category (red color). We quantify more systematically these similarities and differences between barycenters in the following.
        Lastly, the patients groups barycenters exhibit larger fragmentation of actions, with more actions and a larger overall complexity of their barycenters compared to the control group barycenters, which likely reflects significant differences in sequences lengths across groups. We quantify these 

        \paragraph{Vocabulary coverage analysis.\\}
        We next analyze the vocabulary coverage of the estimated barycenters, to assess whether they effectively summarize the diversity of actions performed during the cooking task. We report in \Cref{fig:chapter_6_vocab_missing} the number of missing vocabulary items per type of barycenters, defined as the number of vocabulary items that do not appear in a given barycenter.
        Barycenters exhibited an average of $0.60$ (STD $= 0.87$) missing vocabulary items per estimated barycenter, with 59\% of instances containing the full vocabulary.

        \begin{figure}[!htbp]
            \centering
            \includegraphics[width = 1.0\linewidth]{chapters/chapter 6/figures/chapter_6_vocab_missing.png}    
            \caption{\textbf{Missing vocabulary elements across diagnosis group barycenters.} Each bar represents the proportion of barycenters missing vocabulary items for each diagnosis group. TBI: traumatic brain injury; RIL: radiation-induced leukoencephalopathy.}
            \label{fig:chapter_6_vocab_missing}
        \end{figure}

        For example, the "look kitchen" and "exogeneous" egocentric vision symbols more associated with patients with RIL and less with the control group in \Cref{fig:per_label_pathologie_p_values}, which is effectively reflected in these two type of barycenters. These results indicate that the DBA algorithm produces non-degenerate barycenters that preserve the majority of vocabulary elements across groups.


        \paragraph{Analysis of the composition of barycenters across diagnosis groups.\\}

        Next, we quantify in \Cref{fig:chapter_6_usage_comparison} the relative usage of each label within barycenters by computing the average proportion of time each label appeared across diagnosis group barycenters. Differences across groups were assessed using pairwise Mann-Whitney U tests, correcting multiple comparisons using the Bonferroni-Hockberg procedure, and with effect sizes estimated via rank-biserial correlation. 
        
        \begin{figure}[!htbp]
            \centering
            \includegraphics[width = 1.0\linewidth]{chapters/chapter 6/figures/chapter_6_usage_comparison.png}    
            \caption{\textbf{Comparison of the normalized total duration (that is, proportion of time) of vocabulary items across diagnosis groups barycenters.} Each boxplot shows the distribution of the proportion of barycenters symbols associated with each label across diagnosis groups. Pairwise group comparison uses Mann-Whitney U statistics with Bonferroni-Hockberg corrected p-values. Level of significance is set to $p < 0.05$. TBI: traumatic brain injury; RIL: radiation-induced leukoencephalopathy.}
            \label{fig:chapter_6_usage_comparison}
        \end{figure}
        
        Remarkably, the relative proportion of actions related to "reading the recipe", "looking at the [kitchen/shelf/oven]", "open box ingredients" and "manipulating the oven" were significantly positively associated with commonly acknowledged diagnosis severity, and "using the sink", "merging the butter with the chocolate", "plugging the mixer", and "breaking the chocolate" were negatively associated. Additionally, significant differences between control and patients groups appeared for specific pathology only: for example, difference in "use mixer" and "pour floor/sugar" appears specifically with the TBI barycenters,  while "merge flour/sugar" or "adding the butter" differ significantly with the RIL barycenters only.
        
        Overall, these findings seem to corroborate and complement the findings of \Cref{fig:chapter_6_vocab_missing}, suggesting that our method computes barycenters effectively averaging participant's symbolic representations per diagnosis groups, while helping to quantify characteristic differences associated with each pathology.


        \subsubsection{Comparison of participants symbolic sequences with the diagnosis group barycenters}

        % Behavioral phenotyping via diagnosis group barycenters computation
        We next analyze the TW-TWE distances and proportions of matches between sequences across type of barycenters, to assess quantitatively (\textbf{Q1}) the quality of the averaging procedure, (\textbf{Q2}) the representativity of barycenters within each group, and (\textbf{Q3, Q4}) whether comparisons metrics can be used to differentiate between diagnosis groups. Comparisons of interest are illustrated in \Cref{fig:chapter_6_comparison_barycenter}. We report comparisons of TW-TWE distance and proportion of matching symbols as boxplots in \Cref{fig:chapter_6_classification_report}-(a, b), as well as a statistical analysis in \Cref{tab:mannwhitney_medians_distance} and \Cref{tab:mannwhitney_medians_match_proportions}, respectively. 


        \paragraph{Analysis of the TW-TWE distance.\\}

        While training sequences exhibit lower TW-TWE distance than testing sequences to their group barycenters, these differences were not significant (see lines \#[1-3] in \Cref{tab:mannwhitney_medians_distance}), suggesting the generalizability of the barycenters in the sense of the TW-TWE distance. Additionally, although they did not differ significantly, we note that the distances between patients's sequences and the control barycenters are positively associated with the commonly acknowledged severity of dysexecutive syndromes (line \#6). %CODE/ add something here ? 
        However, we remark that barycenters of patients groups exhibited comparable distance with both patient groups (see lines \#7 and \#8) suggesting that the TW-TWE distance is likely to fail separating between patients groups. 
        
        \begin{table}[htbp]
            \small
            \setlength{\tabcolsep}{4pt}
            \begin{tabular}{r l l l r r r r r r}
            \toprule
            \# & \textbf{Bar.} & \textbf{Group 1} & \textbf{Group 2} & $n_1$ & $n_2$ & $Med_{1}$ & $Med_{2}$ & $U$ & $p$ \\
            \midrule
            1 & Control & Train & Test & 780  & 780  & 1399.25 & 1427.48 & 288754.0  & $8.25\times10^{-2}$ \\
            2 & TBI& Train & Test& 1740 & 1800 & 1731.30 & 1775.63 & 1536885.0 & $3.38\times10^{-1}$ \\
            3 & RIL& Train & Test & 1080 & 1140 & 1739.71 & 1782.22 & 590924.0  & $1.02\times10^{-1}$ \\
            4 & Control & Control$^*$ & TBI & 780 & 3540 & 1427.47 & 1618.69 & 1787872.0 & $3.65\times10^{-38}$ \\
            5 & Control & Control$^*$ & RIL& 780 & 2220 & 1427.47 & 1642.99 & 1131588.0 & $2.34\times10^{-37}$ \\
            6 & Control & TBI & RIL  & 3540 & 2220 & 1618.69 & 1642.99 & 4010360.0 & $1.87\times10^{-1}$ \\
            7 & TBI & TBI$^*$ & RIL& 1800 & 2220 & 1775.63 & 1780.54 & 1988757.0 & $8.01\times10^{-1}$ \\
            8 & TBI  & Train & Control & 1740 & 1560 & 1731.30 & 1595.90 & 1605366.0 & $1.07\times10^{-19}$ \\
            9 & RIL & TBI  & RIL$^*$ & 3540 & 1140 & 1799.25 & 1782.22 & 2081953.0 & $1.06\times10^{-1}$ \\
            10 & RIL & Train  & Control & 1080 & 1560 & 1739.70 & 1657.00 & 938112.0 & $6.68\times10^{-7}$ \\
            \bottomrule
            \end{tabular}
            \caption{\textbf{Mann--Whitney U tests comparing distances between sequences and their corresponding barycenters across training and testing splits, and between diagnosis groups.} The line number ($\#$) facilitates references to individual pairwise comparisons in the text. Bar.: Barycenter; TBI: traumatic brain injury; RIL: radiation-induced leukoencephalopathy. Med: Median, $^*$: comparisons including testing sequences only. }
            \label{tab:mannwhitney_medians_distance}
        \end{table}

        Importantly, while we observe that testing control participants are consistently and significantly closer to their barycenters than both patients groups (lines \#4 and \#5), we observed that control sequences were also found to be consistently and statistically significantly closer to the patient's group barycenters than the training sequences (lines \#8 and \#10). We believe this is due to the "simpler" nature of averaged sequences, which tend to be more similar to the shorter and thus "simpler" sequences of control participants. This influence of sequences length constitutes a limitation that precludes the use of the distance-to-the-barycenter as an effective feature to distinguish sequences of different groups, which needs further investigation in future work.
        
        
        \paragraph{Analysis of the proportion of matching symbols.\\}

        We first note that the training sequences used to form the barycenters exhibited significantly higher proportion of matching symbols compared to the testing sequences (see lines \#[1-3] in \Cref{tab:mannwhitney_medians_match_proportions}). This suggests that barycenters effectively capture the composition of their training sequences, but fail to generalize to the sequences within their diagnosis group. Moreover, similar to the distances, we observed that sequences of control participants consistently exhibited a higher proportion of matching symbols with patient groups' barycenters (lines \#8 and \#10). 

        \begin{table}[htbp]
            \small
            \setlength{\tabcolsep}{4pt}                
            \begin{tabular}{r l l l r r r r r r}
            \toprule
            \# & \textbf{Bar.} & \textbf{Group 1} & \textbf{Group 2} & $n_1$ & $n_2$ & $Med_{1}$ & $Med_{2}$ & $U$ & $p$ \\
            \midrule
            1 & TBI & Train & Test & 1740 & 1800 & 0.538 & 0.482 & 2130390.5 & $6.16\times10^{-77}$ \\
            2 & RIL & Train & Test & 1080 & 1140 & 0.539 & 0.458 & 884591.0  & $5.01\times10^{-71}$ \\
            3 & Control & Train & Test & 780  & 780  & 0.643 & 0.552 & 464016.5  & $3.70\times10^{-72}$ \\
            4 & Control & Control$^*$& TBI & 780  & 3540  & 0.553 & 0.479 & 886225.0  & $2.12\times10^{-55}$ \\
            5 & Control & Control$^*$& RIL & 780  & 2220  & 0.552 & 0.453 & 457685.5  & $1.23\times10^{-85}$ \\
            6 & Control & TBI &  RIL  & 3540 & 2220& 0.479 & 0.453 & 3335527.5 & $4.11\times10^{-22}$ \\
            7 & TBI & TBI $^*$     & RIL & 1800 & 2220 & 0.482 & 0.465 & 2234758.0 & $9.79\times10^{-11}$ \\
            8 & TBI & Train & Control & 1740 & 1560 & 0.538 & 0.512 & 1624527.0 & $1.33\times10^{-22}$ \\
            9 & RIL & RIL$^*$ & TBI & 1140 &  3540 & 0.458  & 0.467 & 2106331.5 & $2.57\times10^{-2}$  \\
            10 & RIL & Train & Control & 1080 & 1560 & 0.539 & 0.486 & 1115676. & $1.03\times10^{-45}$  \\
            \bottomrule
            \end{tabular}
            \caption{\textbf{Mann--Whitney U tests comparing the proportion of matching symbols between sequence-to-barycenter alignments across training and testing splits, and between diagnosis groups.} The line number ($\#$) can be used to refer to individual pairwise comparisons in the text. Bar.: Barycenter; TBI: traumatic brain injury; RIL: radiation-induced leukoencephalopathy. Med: Median, $^*$: testing sequences only. }
            \label{tab:mannwhitney_medians_match_proportions}
        \end{table}

        However, the proportion of matching symbols, appears as a more sensitive feature to capture to group differences. We observe in line \#6 that sequences of patients with TBI exhibited a significantly higher proportion of matching symbols with the control barycenters compared to the RIL participants, and in lines \#7 and \#9 that testing participants of the TBI and RIL groups exhibit larger proportion of matching symbols with their respective barycenters. This suggests that the proportion of matching symbols could be a more effective feature to differentiate between diagnosis groups compared to the TW-TWE distance. We investigate this further below, by evaluating the discriminative power of proportion of matching symbols for pairwise classification of diagnosis groups.  
        
        
        \paragraph{Diagnostic accuracy of the proportion of matching symbols for pairwise classification of diagnosis groups.\\}

        We summarize the classification performance across diagnosis groups in \Cref{fig:chapter_6_classification_report}, showing the ROC curves and AUC values in subfigure (c), confusion matrices in subfigures (d) and (e). We observe that the proportion of matching symbols to the barycenters yields high classification performance for differentiating participants with dysexecutive syndromes from the control participants, with AUC values of $0.77$ (95\% CI: $[0.69, 0.86]$), with higher discriminative power when restricting the analysis to patients with RIL only (AUC=$0.84$ (95\% CI: $[0.75, 0.94]$). Distinguishing between patients with TBI and RIL yielded lower performance, with an AUC of $0.59$ (95\% CI: $[0.48, 0.71]$). These results suggest that the proportion of matching symbols to diagnosis group barycenters effectively captures characteristic differences in task execution between groups, and could serve as a useful feature for behavioral phenotyping.

            
        \begin{figure}[!htbp]
            \centering
            \includegraphics[width = 1.0\linewidth]{chapters/chapter 6/figures/chapter_6_classification_report.png}    
            \caption{\textbf{comparison of TW-TWE distance (a) and proportion of matching symbols (b) between participants symbolic sequences and group-specific barycenters per pathology, (c) receiving operating curves (ROCs) and associated  area under the curves (AUCs) between pairwise diagnostic classifications using the proportion of matching symbols, with associated confusion matrices between coarse (d) and fine-grained (e) predictions.} Prediction of  diagnosis groups are based on the proportion of matching symbols between testing sequences and group-specific barycenters, with AUC 95\% confidence intervals are computed from the Hanley McNeil method. RIL: Radiation-induced Leukoencephalopathy, TBI: Traumatic brain injury.}
            \label{fig:chapter_6_classification_report}
        \end{figure}


\section{Discussions}
    \label{sec:chapter_6_conclusion}

    In this chapter, we presented an adaptation to the TWE distance to quantify the similarity between symbolic sequences of egocentric vision, which relies on the Wasserstein distance between the distribution of symbols occurrences to account for their temporal distribution during the cooking task. Although the distance remains sensitive to sequence length, it yields meaningful optimal alignments that enable comparison based on the amount of matching symbols. These alignments were leveraged in an adaptation of the DBA algorithm to compute average symbolic representations characterizing the typical execution of the task within each diagnosis group. Altogether, these findings demonstrate the feasibility and effectiveness of combining edit distance–based alignment with DBA averaging to derive meaningful average symbolic sequences. The proportion of matching symbols between participant's sequences and diagnosis group barycenters emerged as a discriminative feature to differentiate control participants from patients with dysexecutive syndromes across severity-stratified diagnosis subgroups.
    
    Nevertheless, further work is required to enhance both the robustness and interpretability of the TW-TWE distance and the representativeness of the resulting barycenters. A more systematic analysis of the discrepancies and fine-grained deviations between individual sequences and their group barycenters could ultimately support more objective and scalable behavioral phenotyping of dysexecutive syndromes in ecological contexts.

    Future work could explore the development of more interpretable similarity metrics between barycenters and participant sequences, such as label-wise matching proportions, or identify discriminative subsequences that better capture group-specific execution strategies or DS-induced behavioral alterations.  Refinement of the barycenter computation through neighborhood-aware approaches such as ShapeDBA \cite{fawaz2023shapedba} represents another direction for improvement  Likewise, the robusteness and accuracy of the TW-TWE distance could benefit from the integration of additional edit operations (e.g., transpositions \cite{Damerau1964}) or the adoption of a novel edit distance measure accouting for editions at larger scale than single symbols \cite{Petty2022}, such that behaviors better represented by symbolic subsequences can be compared more effectively.



    %We believe that the superiority of the matching proportion over the distance for classification may stem from the fact that the former directly quantifies the similarity of behaviors in terms of shared actions, while the latter may be more sensitive to sequences interruptions. 


    % Add the Damerau–Levenshtein distance for accoutning in transposition 

    % \cite{10.1007/978-3-031-19778-9_38}: My View is the Best View: Procedure Learning from&nbsp;Egocentric Video

\section{Conclusion}



    Finally, the use of pattern mining techniques could be explored to identify frequently occurring sub-sequences within symbolic representations, which may provide additional insights into common behavioral patterns and strategies employed by participants during the cooking task. These patterns could then be compared across diagnosis groups to identify differences in task execution strategies. Additionally, the integration of other modalities (e.g., audio, eye-tracking, exocentric video) with symbolic representations could be further investigated to enrich the behavioral characterization of participants and provide a more comprehensive understanding of their cognitive and executive functioning. Future work could also explore the use of machine learning learning techniques based on these pseudo-labels to classify participants based on their symbolic representations and identify features that are most predictive of dysexecutive syndromes. Another promising approach may be to leverage the semantic lift enabled by recent advances in language-aware approaches to query or perform language . This is facilitated by the recovered interpretability of symbols semantics through efficent visualizations of independant frames. 

    We believe this framework provides a promising avenue for the analysis of complex neurological disorders having subtle behavioral manifestations, which can ultimately improve our understanding of these conditions and improve patient care.


% \section{Multimodal integration using symbolic representations}
% \label{sec:chapter_6_multimodal_integration}

%     Behavioral expression of the spontaneous and multifaceted nature of DS-related behavioral alterations, and their expression through diverse and semantically rich errors requires a more systematic, holistic approach to multimodal integration.



%     This part of the thesis focuses on developing agnostic computational methods to characterize participant's behavior by integrating the information stemming from complementary modality components. 
%     The first part demonstrated that this integration can successfully be performed by processing each modality first, before combining semantically-meaningful features to obtain behavioral markers. For example, the response to name call variables rely on tailored modality-specific detection algorithms - combining audio and video. Likewise, the social attention or head movements-related variables are obtained from a customized combination of gaze, facial landmarks and head pose information, and aggregated across time window defined by the protocol design.

%     %


%     % From intro of chapter 5: We therefore compute the symbolic representations from these recordings, and leave for future work the addition of exocentric views, which is discussed in \Cref{sec:chapter_5_conclusion}. 
%     In \Cref{sec:chapter_6_multimodal_integration}, we show that our egocentric vision symbolic representation provides sufficiently granular information about the behavioral sequence to serve as a reference modality to facilitate multimodal integration.
    
%     % Two types of analysis are made possible by these symbolic representations: (i) 
%     On the one hand, interpretability of the finite set of atomic actions (e.g., reading the recipe, talking to the examiner, setting-up the oven, etc.) allow to contextualize the analysis of other modality-specific features (such as gaze and oculometric patterns, speech, exocentric motion cues, etc), and allow neurologists to test hypothesis more systematically.

%     On the other hand, other modality-specific features across time — e.g. pupils diameters, gaze patterns, hand ladmarks, and speech patterns, can integrate these symbolic representations as additional and complementary dimensions that enrich the characterization. 

%     Ultimately, we argue that well-resolved, multimodal symbolic representations provide a powerful low-dimensional abstraction of the participant's behavior, unlocking the application of pattern discovery methods to support behavioral phenotyping.




%     Analyze de l'audio: rend possible la l'analyze contextuelle et systematique des interactions sociales, des langages egocentriques (self-directed), et troubles du langages et de la parole (paraphasie, palilalie, mutisme, écholalie, etc), du timbre, et des etats emotionelles. 

%     Analyze multi-vues: rend possible l'étude contextuelle et tri-dimensionnelle des mouvements du corps et des interactions patients-environnements
%     Hands landmarks tracking pour l'analyze de la motricité fine (cf Alexander Krebeich)
%     Skeleteton landmarks tracking pour l'analyze de la motricité grossière

%     Analyze de l'oculomotricité et du regard: rend possible l’étude contextuelle (!) de la charge cognitive, des stratégies attentionnelles, des mécanismes d’anticipation sensorimotrice, et des altérations du contrôle visuo-moteur associées aux pathologies neurologiques, motrices, ou physiologiques

